\chapter{Background teorico}

\section{Il linguaggio Racing Choreographies}

Racing Choreographies è l’estensione di un linguaggio coreografico con l’obiettivo di integrare nel modello globale un costrutto esplicito di competizione tra processi.

Il linguaggio permette di avere:

\begin{itemize}
    \item definizioni di procedure
    \item coreografie composte da istruzioni, chiamate di procedura e costrutti di controllo locali o su esito di race
    \item un insieme di espressioni elementari
    \item comunicazioni
    \item selezioni di etichetta
    \item assegnamenti locali
\end{itemize}

A queste si aggiunge il costrutto di race, che rappresenta l’estensione principale del modello.

La race descrive una competizione tra due processi che tentano di inviare un valore a un processo ricevente. Si formalizzano tre aspetti distinti:

\begin{itemize}
    \item Determinazione del vincitore ossia quale processo riesce a consegnare per primo il proprio valore.
    \item Memorizzazione dell’esito associata a un identificatore di race.
    \item Gestione del perdente il cui valore non viene scartato ma conservato per una successiva sincronizzazione.
\end{itemize}

La semantica operazionale è definita su configurazioni globali che includono:

\begin{itemize}
    \item la coreografia residua;
    \item uno store delle variabili localizzate;
    \item una struttura ausiliaria per la gestione delle race.
\end{itemize}

Le regole operazionali descrivono come le configurazioni evolvono in seguito all’esecuzione di interazioni, condizionali, chiamate ricorsive e race.

Un elemento importante del linguaggio è la presenza di regole che introducono parallelismo in quanto consentono l’esecuzione di azioni indipendenti in ordine non strettamente sequenziale. Questo è formalizzato tramite:

\begin{itemize}
    \item nozioni di lettura e scrittura delle variabili;
    \item condizioni di indipendenza tra azioni;
    \item regole di delay che permettono il riordinamento di istruzioni non interferenti.
\end{itemize}

Il linguaggio introduce inoltre condizioni di Well-formedness. Queste proprietà assicurano che le coreografie siano semanticamente coerenti rispetto al meccanismo di competizione tra cui: uso lineare degli identificatori di race; vincoli strutturali che garantiscono che ogni processo perdente venga correttamente sincronizzato; assenza di configurazioni inconsistenti nella gestione delle race.

Il linguaggio include anche un calcolo target che rappresenta una possibile semantica a livello di processi locali e una definizione di Endpoint Projection (EPP) che associa a ogni processo la parte di comportamento che lo riguarda. La proiezione garantisce che il comportamento locale risultante preservi la struttura della coreografia definita a livello globale.

Nelle sezioni successive verranno analizzate nel dettaglio la sintassi astratta e la semantica operazionale.

\section{Sintassi astratta}

La sintassi astratta descrive la struttura formale dei programmi indipendentemente dalla loro rappresentazione concreta. Essa definisce le categorie sintattiche del linguaggio e le modalità con cui i costrutti possono essere composti.

\subsection*{Programmi}

Un programma è composto da zero o più definizioni di procedure seguite da una coreografia principale:

\[
P ::= D^* \; ; \; C
\]

dove $D^*$ indica una sequenza (eventualmente vuota) di definizioni di procedura e $C$ rappresenta la coreografia principale.

\subsection*{Definizioni di procedura}

Una procedura è identificata da un nome $X$, da una tupla di processi partecipanti e da un corpo espresso come coreografia:

\[
D ::= X(\bar{p}) = C
\]

dove $\bar{p}$ rappresenta la tupla dei processi coinvolti nella procedura.

\subsection*{Coreografie}

Le coreografie descrivono il comportamento globale del sistema e sono definite dalla seguente grammatica:

\[
\begin{array}{rcl}
C &::=& I \; ; \; C \\
  &\mid& \text{if } p.e \text{ then } C \text{ else } C \\
  &\mid& \text{if } s[k] \text{ then } C \text{ else } C \\
  &\mid& X(\bar{p}) \\
  &\mid& 0
\end{array}
\]

Una coreografia può quindi essere:

\begin{itemize}
    \item un’istruzione seguita da un’altra coreografia;
    \item un condizionale basato sulla valutazione locale di un’espressione $p.e$;
    \item un condizionale basato sull’esito di una race identificata da $s[k]$;
    \item una chiamata di procedura;
    \item il termine $0$, che rappresenta la terminazione.
\end{itemize}

\subsection*{Istruzioni}

Le istruzioni costituiscono le azioni elementari della coreografia:

\[
\begin{array}{rcl}
I &::=& p.e \rightarrow q.x \\
  &\mid& p \rightarrow q[l] \\
  &\mid& p.x = e \\
  &\mid& s[k] : p.e \; q.e' \; \perp \; s.x \\
  &\mid& s[k] : p \Rightarrow s.x
\end{array}
\]

Le forme rappresentano rispettivamente:

\begin{itemize}
    \item comunicazione di un valore dal processo $p$ al processo $q$;
    \item selezione di etichetta;
    \item assegnamento locale;
    \item race tra due processi $p$ e $q$ con ricevente $s$;
    \item discharge del processo perdente della race.
\end{itemize}

\subsection*{Espressioni}

Le espressioni sono limitate a valori e variabili:

\[
e ::= v \mid x
\]

dove $v$ rappresenta un valore e $x$ una variabile localizzata.

\subsection*{Metavariabili}

Le metavariabili utilizzate nella definizione sintattica sono:

\begin{itemize}
    \item $p,q,r,s$: processi;
    \item $\bar{p}$: tupla di processi;
    \item $x$: variabile locale a un processo;
    \item $v$: valore;
    \item $l$: etichetta;
    \item $k$: identificatore di race;
    \item $X$: nome di procedura.
\end{itemize}

\section{Semantica operazionale}

Per l’esecuzione del linguaggio è stata definita una semantica operativa semplificata, progettata per essere direttamente implementabile tramite una macchina astratta.

La configurazione globale assume la forma:

\[
C, \Sigma, M
\]

dove:

\begin{itemize}
    \item $C$ è la coreografia residua;
    \item $\Sigma$ è lo store delle variabili localizzate;
    \item $M$ è una memoria esplicita delle race.
\end{itemize}

La relazione di transizione è:

\[
C, \Sigma, M \longrightarrow C', \Sigma', M'
\]

\subsection*{Race}

\textbf{Regola $\langle l\text{race} \rangle$}

Premesse:
\[
s[k] \notin dom(M)
\]
\[
\Sigma(p,e) \downarrow v
\]
\[
\Sigma(q,e') \downarrow v'
\]

Transizione:
\[
C, \Sigma, M, \; s[k] : p.e \; q.e' \perp s.d ; C
\longrightarrow
C, \Sigma[s.d \mapsto v], M[s[k] \mapsto (L,R,v,v')]
\]

Se la race identificata da $s[k]$ non è ancora presente in $M$ e le espressioni vengono valutate rispettivamente a $v$ e $v'$, allora nel caso di vittoria del lato sinistro viene assegnato a $s.d$ il valore $v$. In $M$ viene memorizzata una tupla contenente l’informazione sul vincitore ($L$), sul perdente ($R$), sul valore vincente e sul valore perdente.

\textbf{Regola $\langle r\text{race} \rangle$}

Premesse: 
\[
s[k] \notin dom(M)
\]
\[
\Sigma(p,e) \downarrow v
\]
\[
\Sigma(q,e') \downarrow v'
\]

\[
C, \Sigma, M, \; s[k] : p.e \; q.e' \perp s.d ; C
\longrightarrow
C, \Sigma[s.d \mapsto v'], M[s[k] \mapsto (R,L,v',v)]
\]

Nel caso di vittoria del lato destro, viene assegnato a $s.d$ il valore $v'$ e in $M$ viene registrata la tupla corrispondente con indicazione del lato vincitore e dei valori coinvolti.

\subsection*{Comunicazione}

\textbf{Regola $\langle com \rangle$}

Premessa:
\[
\Sigma(p,e) \downarrow v
\]

Transizione:
\[
C, \Sigma, M, \; p.e \rightarrow q.x ; C
\longrightarrow
C, \Sigma[q.x \mapsto v], M
\]

Il processo $p$ valuta l’espressione $e$ ottenendo $v$. Il valore viene assegnato alla variabile $q.x$ nello store. La memoria $M$ non viene modificata.

\subsection*{Selezione}

\textbf{Regola $\langle sel \rangle$}

\[
C, \Sigma, M, \; p \rightarrow q[l] ; C
\longrightarrow
C, \Sigma, M
\]

La selezione di etichetta non modifica né lo store né la memoria delle race; l’esecuzione prosegue con la coreografia successiva.

\subsection*{Assegnamento}

\textbf{Regola $\langle asg \rangle$}

Premessa:
\[
\Sigma(p,e) \downarrow v
\]

Transizione:
\[
C, \Sigma, M, \; p.x = e ; C
\longrightarrow
C, \Sigma[p.x \mapsto v], M
\]

Il valore dell’espressione viene assegnato alla variabile locale $p.x$ nello store.

\subsection*{Condizionali locali}

\textbf{Regola $\langle cndt \rangle$}

Premessa:
\[
\Sigma(p,e) \downarrow true
\]

Transizione:
\[
C, \Sigma, M, \; \text{if } p.e \text{ then } C_1 \text{ else } C_2
\longrightarrow
C_1, \Sigma, M
\]

Se l’espressione valutata localmente è vera, viene selezionato il ramo then.

\textbf{Regola $\langle cndf \rangle$}

Premessa:
\[
\Sigma(p,e) \downarrow false
\]

Transizione:
\[
C, \Sigma, M, \; \text{if } p.e \text{ then } C_1 \text{ else } C_2
\longrightarrow
C_2, \Sigma, M
\]

Se l’espressione è falsa, viene selezionato il ramo else.

\subsection*{Condizionali su race}

\textbf{Regola $\langle rcndt \rangle$}

Premesse:
\[
s[k] \in dom(M)
\]
\[
M[s[k]] = (L, \_, \_, \_)
\]

Transizione:
\[
C, \Sigma, M, \; \text{if } s[k] \text{ then } C_1 \text{ else } C_2
\longrightarrow
C_1, \Sigma, M
\]

Se la memoria indica che il lato sinistro ha vinto, viene eseguito il ramo then.

\textbf{Regola $\langle rcndf \rangle$}

Premesse:
\[
s[k] \in dom(M)
\]
\[
M[s[k]] = (R, \_, \_, \_)
\]

Transizione:
\[
C, \Sigma, M, \; \text{if } s[k] \text{ then } C_1 \text{ else } C_2
\longrightarrow
C_2, \Sigma, M
\]

Se la memoria indica che ha vinto il lato destro, viene eseguito il ramo else.

\subsection*{Discharge}

\textbf{Regola $\langle dis \rangle$}

Premesse:
\[
s[k] \in dom(M)
\]
\[
M[s[k]] = (w, \ell, v_{win}, v_{lose})
\]

Transizione:
\[
C, \Sigma, M, \; \text{discharge } s[k] : \ell \rightarrow s.x ; C
\longrightarrow
C, \Sigma[s.x \mapsto v_{lose}], M
\]

Il valore del processo perdente, memorizzato in $M$, viene assegnato alla variabile $s.x$. La memoria $M$ non viene rimossa, ma può essere marcata come già utilizzata a livello implementativo.

\subsection*{Chiamata di procedura}

\textbf{Regola $\langle call \rangle$}

Premessa:
\[
\Delta(X) = (\bar{p}_{form}, C_{body})
\]

Transizione:
\[
C, \Sigma, M, \; \text{call } X(\bar{p}_{act}) ; C
\longrightarrow
C, \Sigma, M, \; subst(C_{body}, \bar{p}_{form} \mapsto \bar{p}_{act})
\]

La chiamata di procedura viene realizzata tramite sostituzione dei parametri formali con quelli attuali nel corpo della procedura.

\subsection*{Terminazione}

\textbf{Regola $\langle end \rangle$}

\[
C, \Sigma, M, \; 0
\longrightarrow
C, \Sigma, M, \; 0
\]

La configurazione terminale non produce ulteriori transizioni.